% ==============================================================================
% secoes/03_metodologia.tex - SEÇÃO III: METODOLOGIA
% Tema: Federico Faggin e o Primeiro Microprocessador Monolítico (Intel 4004)
% ==============================================================================

\section{Metodologia}
\label{sec:metodologia}

\subsection{Tipo de Estudo}
\label{subsec:metodologia_abordagem}

Trata-se de pesquisa histórica e analítica, de base técnica e documental, voltada à interface entre a física dos dispositivos semicondutores e a arquitetura de computadores.

\subsection{Fontes Primárias e Secundárias}
\label{subsec:metodologia_coleta}

As fontes primárias consistem em:
\begin{itemize}
\item patentes industriais de Federico Faggin sobre dispositivos MOS e métodos de fabricação com porta de silício, como a US Patent 3.475.234 \cite{faggin1969patent}, e a patente US 3.821.715 sobre o sistema de memória do computador multichip \cite{faggin1974patent};
\item manuais técnicos e notas de aplicação da Intel, em particular o \textit{MCS-4 Micro Computer Set Users Manual} \cite{intel1973mcs4};
\item documentos autobiográficos e artigos memoriais técnicos de Faggin sobre o desenvolvimento da arquitetura e das regras de layout do 4004 \cite{faggin1996, faggin2021}.
\end{itemize}

As fontes secundárias compreendem:
\begin{itemize}
\item periódicos especializados em história da tecnologia e dos semicondutores, como \textit{IEEE Micro} e \textit{Proceedings of the IEEE} \cite{noyce1981, cavin2012};
\item tratados clássicos de circuitos integrados e arquitetura de computadores \cite{mead1980, weste2010};
\item estudos socioeconômicos sobre automação industrial e difusão de tecnologias digitais nas últimas décadas do século XX.
\end{itemize}

\subsection{Procedimento de Análise}
\label{subsec:metodologia_triangulacao}

O procedimento consiste na triangulação sistemática entre três conjuntos de dados: (i) a física de dispositivos semicondutores; (ii) os parâmetros analíticos dos circuitos integrados, extraídos de patentes e manuais; e (iii) os relatórios históricos sobre a difusão industrial dos microprocessadores. As relações analíticas da \cref{sec:resultados} são tratadas como modelos de primeira ordem, usados para explicitar tendências, e não para reproduzir valores de projeto.
